\section{Case Studies: Extreme Atmospheric Events}
\label{sec:case_studies}
To further demonstrate the practical utility of PhyST-Net, we analyze its performance during several representative high-impact events. These case studies highlight the model's ability to maintain grid stability where traditional data-driven models often fail due to physical boundary violations.

\subsection{Case Study A: Rapid Wind Ramp at Hill of Towie}
\label{sec:case_ramp}
Wind ramps are defined as sudden shifts in wind velocity that can exceed 50\% of rated capacity within minutes. During a recorded low-pressure front at the Hill of Towie (HoT) site, wind speeds jumped from 4 m/s to 18 m/s in less than 20 minutes. Standard Transformer models exhibited significant ``Overshoot,'' predicting power outputs 12\% above the rated mechanical limit. This was caused by the model's attempt to follow the steep gradient without awareness of the turbine's pitch control saturation.

In contrast, PhyST-Net's AGSP module adaptively reduced its receptive field to capture the rapid acceleration, while the PI-DCH projected the output strictly onto the $P_{rated}$ boundary. The resulting forecast was physically consistent and provided grid operators with a reliable signal for fast-frequency response (FFR) activation.

\subsection{Case Study B: Large-Scale Cluster Synchronization in SDWPF}
\label{sec:case_cluster}
In massive farms like SDWPF (134 turbines), a weather front moving across the site affects different turbine clusters at different times. Traditional models with static graphs often ``average out'' these temporal shifts, leading to a blurry farm-wide forecast that misses the peak ramp timing. PhyST-Net's R-STMF Trend branch correctly identified the propagation delay across the 134 turbines, learning to ``anticipate'' the power increase in Eastern clusters based on observations from Western clusters. This ``look-ahead'' capability improved the farm-wide peak power timing accuracy by 25\%.

\subsection{Case Study C: NWP Sensor Bias Correction in Kelmarsh}
\label{sec:case_nwp}
During a regional storm event at Kelmarsh, the forecasted wind speed from the numerical weather prediction (NWP) model was consistently 2.5 m/s higher than the actual recorded anemometer data due to local topographical shadowing. Standard MLP-based modulation attempted to scale the entire hidden state based on this biased feature, leading to a systematic over-prediction. 

The KAN-enhanced K-AMC layer, however, utilized its learnable B-splines to identify the specific wind regime where the NWP was unreliable. By applying a non-linear ``dampening'' to the weather features in that specific speed range, PhyST-Net reduced the forecasting bias by 60\% compared to standard MLP-based modulation. This demonstrates that Kolmogorov-Arnold networks (KANs) can learn to correct sensor inaccuracies through local spline adjustments without affecting the model's general performance.

\subsection{Case Study D: Low-Wind Persistence and Negative Power Avoidance}
\label{sec:case_lowwind}
During a period of atmospheric stagnation at the SDWPF site, wind speeds dropped below the 3 m/s cut-in threshold for several hours. Standard data-driven models, optimized solely on MSE, produced oscillating predictions that included negative power values of up to -2\% of rated capacity. While mathematically small, these values are physically non-sensical and can cause numerical errors in downstream grid optimization software.

PhyST-Net's physics-informed dynamic capping head (PI-DCH) effectively mapped these latent oscillations to a strict zero-output state. By maintaining the prediction manifold within the [0, $P_{rated}$] domain, the model provided a clean, physically plausible signal that required no post-processing by the grid operator. This highlights the value of structural physics awareness in safety-critical applications.
